\section{INTRODUCTION}
\IEEEPARstart{C}{yber-physical} systems (CPS) couple sensing, communication, computation, and control with physical processes. After a cross-layer attack, recovery controllers must act before all compromised components can be isolated. A difficult case arises when the recovery mechanism itself depends on distributed learning: local observations may be corrupted, communication may be delayed, and the federated knowledge used to update recovery policies may be poisoned by unreliable or Byzantine participants. In this setting, recovery is no longer only a control-performance problem. It becomes a \emph{recovery-integrity} problem: the controller must restore the plant while preventing untrusted learned knowledge from propagating into unsafe physical actions.

Existing reinforcement-learning-based recovery methods are commonly centralized or assume that the policy-training pipeline remains trustworthy after the attack. Federated reinforcement learning (FRL) reduces raw-data exchange, but standard FRL introduces a new attack surface. A malicious participant can submit arbitrary model updates or metadata, heterogeneous clients can generate incompatible updates, and a server may aggregate stale or poisoned knowledge into a global policy. Data locality does not prevent this failure mode. Likewise, empirical robustness of an aggregation rule does not by itself imply that the resulting control action is safe.

This observation motivates a different design principle: \emph{learning should be treated as untrusted, whereas execution must remain certifiably safe}. We formulate trustworthy heterogeneous federated recovery, where subsystem-specific observations, adapters, recovery updates, and client-side confidence signals are not assumed to be globally trustworthy. The shared recovery model is used to improve recovery utility, but every learned action is constrained by a physical safety-recovery contract before execution. Consequently, Byzantine robustness protects learning quality, whereas a safety certificate protects the physical plant even when the learned policy is poor or adversarially poisoned.

The revised framework retains heterogeneous local observation and action adapters, because different CPS subsystems need not expose compatible state or actuation dimensions. Only a shared recovery backbone is federated. Recovery-state alignment remains useful for reducing honest-client drift, but prototype similarity and client-reported advantage are no longer treated as security roots. Instead, honest clients use recovery certificates to decide which local knowledge is eligible for sharing, while the server applies Byzantine-resilient aggregation under an explicit adversarial budget. This removes the circular assumption that credibility-weighting bias is bounded simply because the credibility mechanism is expected to work.

A second issue is compression. Communication-efficient FRL often uses Top-$k$ sparsification, yet prior analyses commonly introduce an abstract contraction coefficient. For deterministic Top-$k$, the contraction need not remain an unverifiable assumption: retaining the $k$ largest-magnitude coordinates of a $d$-dimensional vector guarantees a residual-energy bound with $\delta_k=k/d$. We therefore keep the compression residual explicit and combine it with a mechanism-derived Byzantine aggregation deviation rather than absorbing both effects into generic stochastic-gradient noise.

The resulting architecture separates three trust layers. At the \emph{local learning layer}, subsystem-specific adapters and recovery-value estimation produce a candidate update. At the \emph{federation layer}, only certificate-gated shared-backbone updates are communicated and a robust aggregator limits Byzantine influence. At the \emph{execution layer}, a robust safety-recovery shield projects the learned action onto a feasible set defined by physical safety and recovery constraints. This separation gives the paper a narrower but stronger goal: we do not claim that federated learning itself is perfectly secure or private; we guarantee that an unsafe learned proposal cannot directly violate the modeled safety contract when the certificate assumptions hold.

The present work makes four contributions.
\begin{itemize}
    \item We formulate \emph{trustworthy heterogeneous federated recovery} for CPS under cross-layer attacks, explicitly modeling both physical-layer corruption and federated-layer Byzantine participants. The formulation distinguishes learning utility, aggregation integrity, and physical safety instead of conflating them in one empirical robustness score.
    \item We redesign AEFC into a certificate-gated heterogeneous federation mechanism. Local observation/action adapters remain private, recovery-state alignment controls honest drift, and only recovery-supported backbone updates are eligible for communication. Server aggregation is Byzantine-resilient and does not trust client-reported credibility, advantage, or prototype similarity as an authentication signal.
    \item We provide mechanism-level guarantees. A robust safety-recovery shield yields policy-independent forward invariance of the safe set under bounded state/attack uncertainty and feasibility. For deterministic Top-$k$, we derive the explicit contraction factor $\delta_k=k/d$; for robust aggregation, we bound the deviation from the honest-update center under an explicit Byzantine budget.
    \item We revise the experimental claims to match the available evidence. IEEE 39-bus and IEEE 118-bus results are used for recovery performance, communication efficiency, heterogeneity stress, and empirical attack robustness. We additionally define the certificate, adaptive-Byzantine, and privacy evaluation protocol required to validate the new trustworthy-recovery guarantees without fabricating unsupported numerical results.
\end{itemize}

The remainder of the paper is organized as follows. Section II reviews resilient CPS recovery, federated and multi-agent reinforcement learning, and trustworthy federated learning. Section III presents the trustworthy heterogeneous recovery model, adversary model, certificate-gated AEFC mechanism, robust aggregation, safety-recovery shield, and theoretical analysis. Section IV reports the supported experiments and the revised security-validation protocol. Section V concludes the paper.
